\section{Regular languages: effective equivalences, closure, and decisions}
\label{sec:regular}

All alphabets in this section are finite, with letters numbered below a supplied
integer $k\geq 1$.  A state set is $[n]=\{0,\ldots,n-1\}$, with $n\geq1$.
A deterministic automaton is a finite transition table
$\delta:[n]\times[k]\to[n]$, an initial state $q_0<n$, and an accepting
bit-vector $F:[n]\to\{0,1\}$.  A nondeterministic automaton replaces
$\delta$ by a Boolean relation $E(q,a,r)$ and permits a bit-vector $I$ of
initial states.  An $\eps$-automaton also has a Boolean relation $E_\eps(q,r)$.
All these data are number codes.  Well-formedness is a bounded test; malformed
codes can be sent to a fixed rejecting automaton.  Consequently every
construction below denotes a total primitive recursive function on codes.

A path is a finite list of states and edge labels.  Its validity, endpoints,
and the word obtained by deleting $\eps$ labels are primitive recursive.
The accepted language is written $L(M)$.  This notation abbreviates an
arithmetic predicate, not an additional sort of arbitrary sets.  In
particular, the equality $L(M)=L(N)$ abbreviates equality of their accepting
bits for every word.  All numerical bounds below are deliberately elementary;
optimal running times are not needed for their interpretation in \ipra.
The corresponding standard Lean interfaces are documented in
\cite{MathlibDFA,MathlibNFA,MathlibEpsilon}; our number-coded implementations
must additionally establish the primitive-recursive and object-theoretic
bridges described below.

\subsection{Finite runs and elimination of nondeterminism}

\begin{lemma}[Run composition]\label{reg:run}
There is a primitive recursive function $\delta^*(q,w)$ satisfying
\[
 \delta^*(q,\eps)=q,\qquad
 \delta^*(q,wa)=\delta(\delta^*(q,w),a),\qquad
 \delta^*(q,uv)=\delta^*(\delta^*(q,u),v).
\]
Thus membership in $L(M)$ is primitive recursive, uniformly in a finite
well-formed DFA code $M$.
\end{lemma}
\begin{proof}
Store the current state and scan the coded input from left to right, taking
one table lookup per letter.  After $i$ letters the stored state is the state
at position $i$ of the unique valid run.  This is an equality between
primitive recursive terms.  Induction on the number of scanned letters
proves the invariant.  Starting the second scan at the state reached by the
first proves the third equation by the same quantifier-free induction.
Acceptance is the final lookup $F(\delta^*(q_0,w))$.
\end{proof}

\paragraph{Finite-graph facts used below.}
The bounded reachability result~\ref{base:graph} applies to every graph
whose vertices are the numbered states of an automaton.  It supplies a
primitive recursive path-shortening operation, a reachability test, and
shortest path witnesses of length below the number of states.  Edge
labels can be retained during each deletion, so a shortened automaton
path gives an explicitly coded shorter word.  The deletion preserves
endpoints, not generally the word itself; where the word must remain
unchanged, we shorten only blocks of $\eps$ edges.

\begin{proposition}[$\eps$-closure has a finite bound]\label{reg:epsilon}
For an $n$-state $\eps$-automaton, the $\eps$-closure $C(S)$ of a state
bit-vector $S$ is primitive recursive.  It is computed by $n$ iterations of
\[
 S\longmapsto S\cup\{r:\exists q<n\;(q\in S\wedge E_\eps(q,r))\}.
\]
It satisfies $S\subseteq C(S)$, $C(C(S))=C(S)$, and preserves finite unions.
If a word of length $\ell$ is accepted, it has an accepting path with at most
$\ell+(\ell+1)(n-1)$ edges.
\end{proposition}
\begin{proof}
After $j$ iterations the bit at $r$ is set exactly when there is an
$\eps$-path of length at most $j$ from $S$ to $r$.  The two directions are
proved by appending an edge or deleting the last edge of a nonempty path.
By the bounded reachability result~\ref{base:graph}, $n-1$ edges already suffice; using $n$ iterations
avoids an exceptional subtraction.  Closure and idempotence follow by
concatenating and shortening paths.  A path from a union starts in one of
its finitely many constituent sets, which proves the union equation.
An accepting path for a length-$\ell$ word has exactly $\ell$ letter edges
and $\ell+1$ intervening $\eps$ blocks.  Shorten each block separately to
length at most $n-1$.  The word and the state at every letter edge are
unchanged, giving the stated bound.  This also replaces the apparently
unbounded existential path definition of acceptance by bounded search.
\end{proof}

\begin{theorem}[Subset construction]\label{reg:subset}
Every $n$-state $\eps$-NFA has an equivalent DFA with $2^n$ states.  The
construction and its correctness proof are primitive recursive and
formalizable in \ipra.  Ordinary NFAs and DFAs consequently define the same
languages as $\eps$-NFAs.
\end{theorem}
\begin{proof}
A DFA state is an $n$-bit subset.  Its initial subset is $C(I)$ and its
letter transition is
\[
 T(S,a)=C\bigl(\{r:\exists q<n\;(q\in S\wedge E(q,a,r))\}\bigr).
\]
A subset accepts when it meets $F$.  For every prefix $u$ of the input,
maintain the invariant that the current bit-vector consists exactly of the
endpoints of paths labelled $u$, allowing a final $\eps$ block.  For the
empty prefix this is the definition of $C(I)$.  For $ua$, split an accepting
path immediately before its last letter edge; conversely concatenate a path
for $u$, that edge, and a final $\eps$ block.  The bounded witnesses from
Proposition~\ref{reg:epsilon} make the invariant a bounded Boolean
assertion for the particular input.  Quantifier-free induction proves it.
A DFA is an NFA with singleton transitions and initial set, and an NFA is
an $\eps$-NFA with no $\eps$ edges.  These embeddings preserve paths.
\end{proof}

\begin{proposition}[Completion and finite languages]\label{reg:finite}
A partial $n$-state DFA can be completed with at most $n+1$ states.  Every
explicit finite list of words $w_1,\ldots,w_t$ defines a regular language,
with a complete DFA using at most $2+\sum_i|w_i|$ states.
\end{proposition}
\begin{proof}
For completion, add a nonaccepting sink, send every missing transition to
it, and loop there on every letter.  A formerly defined run is unchanged;
a run that encounters a missing transition can never accept.
For the second assertion, make states for all distinct prefixes of listed
words, including $\eps$, and add a sink.  From prefix $u$ on letter $a$, go
to $ua$ if it is a listed prefix, otherwise to the sink.  Mark exactly the
listed complete words as accepting.  Induction along a word says that its
run is either at that prefix or at the sink.  There are at most
$1+\sum_i|w_i|$ prefixes; comparing the finitely many codes removes
duplicates constructively.  The empty list and lists containing $\eps$
are handled by the same construction.
\end{proof}

\subsection{Closure constructions}

\begin{theorem}[Boolean closure]\label{reg:boolean}
Regular languages are closed under union, intersection, difference,
symmetric difference, and complement.  Binary operations on $n$- and
$m$-state complete DFAs use $nm$ states; complement uses $n$ states.
\end{theorem}
\begin{proof}
The product automaton has initial state $(q_0,r_0)$ and transition
$(q,r)\mapsto(\delta_M(q,a),\delta_N(r,a))$.  Its run on $w$ is the pair
of the two runs by Lemma~\ref{reg:run}.  For union, intersection,
difference, or symmetric difference, apply the corresponding Boolean
operation to their accepting bits.  Complement replaces $F(q)$ by
$1-F(q)$ and requires no other change.  The complement statement uses
only decidability of a finite accepting bit, not excluded middle for an
arbitrary arithmetic formula.
\end{proof}

\begin{theorem}[Concatenation and star]\label{reg:concatstar}
If $M,N$ have $n,m$ states, there are $\eps$-NFAs for
$L(M)\cup L(N)$, $L(M)L(N)$, and $L(M)^*$ with respectively
$n+m$, $n+m$, and $n+1$ states.  Thus all three languages are regular.
\end{theorem}
\begin{proof}
First rename states so the two copies are disjoint.  For union take the
union of initial sets and the union of final sets, with no edges between
the copies.  Every path stays in its initial copy.
For concatenation use the first initial set, the second final set, and add
$\eps$ edges from every final state of the first copy to every initial
state of the second.  Every accepting path crosses once from the first
copy to the second.  Splitting there gives a factorization $uv$ with
$u\in L(M)$ and $v\in L(N)$; concatenating two witnessing paths gives the
converse.  Empty factors cause no difficulty.
For star add a new state $z$ that is both the sole initial state and the
sole final state.  Add $\eps$ edges from $z$ to every old initial state
and from every old final state to $z$.  A nonempty accepting path splits
at its visits to $z$ into finitely many old accepting paths.  The empty
path represents the zero-factor product.  Conversely concatenate any
finite list of old accepting paths with the connecting $\eps$ edges.
All splits and concatenations act on supplied finite certificates.
Theorem~\ref{reg:subset} then supplies complete DFAs.
\end{proof}

\begin{theorem}[Reversal]\label{reg:reverse}
Regular languages are closed under word reversal.  An $n$-state NFA for
$L$ yields an $n$-state NFA for $L^R$.
\end{theorem}
\begin{proof}
Reverse every edge, interchange the initial and final bit-vectors, and
leave each edge label unchanged.  Reversing a path reverses its sequence
of nonempty labels, and doing it twice returns the original path.
This is a primitive recursive bijection of accepting path certificates.
\end{proof}

\begin{theorem}[Homomorphisms and inverse homomorphisms]\label{reg:morphism}
Let $h$ be specified by a finite table of output words $h(a)$.  Regular
languages are closed under $h$ and $h^{-1}$, including erasing
homomorphisms.  A direct-image construction from an $n$-state DFA uses at
most $n(1+\sum_a|h(a)|)$ $\eps$-NFA states; an inverse-image construction
uses the original number of DFA states.
\end{theorem}
\begin{proof}
For direct image replace every edge labelled $a$ by a private path labelled
$h(a)$.  If $h(a)=\eps$, use an $\eps$ edge; if its length is $d>0$, use
$d$ edges and $d-1$ fresh intermediate states.  Intermediate states belong
to a single replaced edge and have no additional incident edges.  Every
accepting path therefore splits, at the original states, into completed
replacements of original edges.  Its label is $h(u)$ for an accepted
original word $u$.  Conversely replace the edges of any accepting path.
The exact additional state count is
$n\sum_a\max(|h(a)|-1,0)$, bounded as stated.
For inverse image replace the transition on input letter $a$ by
$q\mapsto\delta^*(q,h(a))$.  Run composition gives
$\delta_{\mathrm{new}}^*(q,w)=\delta^*(q,h(w))$, by induction on the
number of letters of $w$.  Initial and final states are unchanged.
\end{proof}

\begin{theorem}[Regular substitution]\label{reg:substitution}
Suppose each source letter $a$ is assigned a regular language $K_a$ by a
finite automaton.  Replacing each letter of words in a regular language
$L$ by an independently chosen word of $K_a$ produces a regular language.
If the source DFA has $n$ states and the automata for the $K_a$ have
$n_a$ states, $n+n\sum_a n_a$ states suffice for an $\eps$-NFA.
\end{theorem}
\begin{proof}
For each source transition $q\mathrel{\smash{\xrightarrow{a}}}r$ take a
private copy of the automaton for $K_a$.  Remove the source edge, put
$\eps$ edges from $q$ to the copy's initial states, and from its accepting
states to $r$.  Keep the source initial and final states.  An accepting
path splits at successive source states into accepting paths through the
private copies, giving a source word and one replacement word per letter.
Conversely, given these finite witnesses, concatenate the corresponding
paths.  A copy can accept $\eps$, but a completed traversal still records
one source letter; no minimum length assumption is used.  There is one
copy per source state and source letter, proving the bound.
\end{proof}

\begin{theorem}[Quotients by a regular language]\label{reg:quotient}
For regular $L,R$, the languages
\[
 L/R=\{u:\exists v\;(v\in R\wedge uv\in L)\},\qquad
 R\backslash L=\{v:\exists u\;(u\in R\wedge uv\in L)\}
\]
are regular.  Given DFAs of $n,m$ states, the existential witnesses needed
to compute the new state sets may always be chosen with length below $nm$.
\end{theorem}
\begin{proof}
For $L/R$, retain the transition table of $L$ and mark $q$ accepting
exactly when a final pair is reachable from $(q,r_0)$ in the product of
$L$ and $R$.  Such a path has length below $nm$ by
the bounded reachability result~\ref{base:graph}.  Therefore its existence is a bounded test, and the
new accepting condition is exactly the displayed definition.
For $R\backslash L$, mark $q$ initial in an NFA exactly when
$(q,r)$ is reachable from $(q_0,r_0)$ for some accepting state $r$ of $R$.
Keep the old transitions and final states of $L$.  The same bound provides
a word $u$ reaching each marked $q$.  A run starting at $q$ on $v$ is
accepting exactly when $uv\in L$.  The two path concatenations prove both
inclusions.  The assertion is intentionally restricted to supplied
regular descriptions of $R$: an arbitrary predicate for $R$ does not
supply a decision procedure for these new finite state sets.
\end{proof}

\begin{corollary}[Prefixes, suffixes, and factors]\label{reg:factors}
The prefixes, suffixes, and contiguous factors of a regular language are
regular.  Their witnessing extension words can be chosen with length
below $n$ on each side, for an $n$-state DFA.
\end{corollary}
\begin{proof}
For prefixes, mark as accepting all states from which a final state is
reachable.  For suffixes, use all states reachable from the original
initial state as NFA initial states.  For factors, make both changes.
A path witnessing an extension can be shortened without changing its
endpoints, so each extension has length below $n$.  Concatenation of the
left extension, the proposed factor, and the right extension proves the
three characterizations.
\end{proof}

\subsection{Regular expressions and regular grammars}

A regular expression is a finite syntax tree generated by
$0,1,a,+,\cdot,{}^*$.  Its denotation interprets these symbols as the empty
language, $\{\eps\}$, a singleton letter, union, concatenation, and star.
A membership certificate records choices at union nodes, word splits at
product nodes, and a finite factor list at star nodes.  Empty factors can
be deleted at a star without changing its product.  For a nonempty word
of length $\ell$, at most $\ell$ remaining factors are needed.  For the
empty word, the zero-factor certificate always suffices.

\begin{proposition}[Bounded regular-expression membership]\label{reg:regexdec}
Membership in a regular expression is primitive recursive uniformly in
the expression and word codes.  It has an algorithm using only a finite
table indexed by expression nodes and input intervals.
\end{proposition}
\begin{proof}
For a fixed word $w$ of length $\ell$, store a Boolean value for every
subexpression $P$ and interval $[i,j)$, $0\leq i\leq j\leq\ell$.
The values for $0,1,a$ are immediate.  Union takes the disjunction.
Product searches the $j-i+1$ splits $i\leq t\leq j$.  For $P^*$ first
set every empty interval true; a nonempty interval is true exactly when
some strictly increasing sequence
$i=t_0<t_1<\cdots<t_r=j$ has every $P$ interval true.  There are at most
$2^{j-i-1}$ such cut sets.  Removing empty factors proves that this
bounded condition is equivalent to star membership.  Evaluate nodes in
postorder, so their children have already been evaluated.  The table has
at most $s(\ell+1)^2$ entries when there are $s$ syntax nodes.  Certificate
extraction stores the first successful split or cut set.  Conversely,
each stored choice assembles a valid certificate.  This proves soundness
and completeness by a bounded conjunction over the table, not by an
induction whose hypothesis quantifies over arbitrary languages or words.
\end{proof}

\begin{theorem}[Regular expressions to automata]\label{reg:thompson}
An expression with $s$ syntax nodes has an equivalent $\eps$-NFA with at
most $2s$ states.  The conversion is primitive recursive.
\end{theorem}
\begin{proof}
Give every subexpression a fresh entry and exit.  Its entry has no incoming
edges inside that subautomaton, and its exit has no outgoing edges there.
For $0$ put no edge between them, for $1$ put an $\eps$ edge, and for a
letter put that labelled edge.  For $P+Q$, connect the new entry to the
two child entries, and the two child exits to the new exit, using $\eps$
edges.  For $PQ$, connect the new entry to $P$'s entry, $P$'s exit to
$Q$'s entry, and $Q$'s exit to the new exit.  For $P^*$ connect the new
entry to the new exit and to $P$'s entry, and $P$'s exit to $P$'s entry
and the new exit.  All these connecting edges are $\eps$ edges.  These
instructions allocate precisely two states per
syntax node.
For union, an entry-to-exit path chooses one child.  For product, it
visits the two child subautomata in order.  For star, split it at returns
to the child entry after the child exit: this gives finitely many child
accepting paths, or the direct zero-factor path.  The converse constructs
paths from membership certificates with the same connectors.  The
entry/exit invariant prevents any other way of moving between children.
For formalization, fix the input word and prove these statements for all
subexpressions and input intervals in the finite table of
Proposition~\ref{reg:regexdec}; shorten $\eps$ blocks when necessary.
Thus the structural proof has an explicitly bounded quantifier-free
invariant.
\end{proof}

\begin{theorem}[Automata to regular expressions]\label{reg:eliminate}
Every finite automaton has an equivalent regular expression, obtained by
a primitive recursive construction.  Together with
Theorem~\ref{reg:thompson}, this proves Kleene's equivalence in \ipra.
\end{theorem}
\begin{proof}
It suffices to treat an $n$-state automaton whose edges carry letters or
$\eps$.  For $i,j<n$ let $R^{(0)}_{ij}$ be the sum of all labels of direct
edges from $i$ to $j$, adding $1$ when $i=j$ and writing $0$ for an empty
sum.  For $0\leq h<n$ define
\[
 R^{(h+1)}_{ij}
 =R^{(h)}_{ij}
  +R^{(h)}_{ih}\bigl(R^{(h)}_{hh}\bigr)^*R^{(h)}_{hj}.
\]
The invariant says that $R^{(h)}_{ij}$ describes exactly the labels of
paths from $i$ to $j$ all of whose internal vertices are below $h$.
At $h=0$ such a path is a single edge or the empty path.  For the induction
step, a path with no internal occurrence of $h$ uses the first summand.
Otherwise split at its first and last internal occurrences of $h$ and at
all the occurrences in between.  The first and last pieces have no
internal $h$, and the middle pieces are $h$-to-$h$ paths with no internal
$h$.  Their labels belong to the three factors of the second summand.
Conversely concatenate these paths; every new internal vertex is either
an old one below $h$ or $h$ itself.  This argument permits $i=h$, $j=h$,
and empty pieces.
At $h=n$ every internal vertex is allowed.  Sum $R^{(n)}_{ij}$ over
initial $i$ and final $j$.  Path splitting and concatenation prove equality
with the accepted language.
The recursion has exactly $n$ rounds and $n^2$ expressions per round.
If the largest base expression has size $b$, the recurrence for sizes is
$s_{h+1}\leq4s_h+4$, hence $s_h\leq4^h(b+2)$.  The final finite sum has
size at most $2n^2(s_n+1)$.  These coarse bounds show directly that even
an implementation that copies syntax trees is primitive recursive.
For each fixed input word the invariant can be expressed over its
finitely many intervals and endpoint pairs; a syntactic path certificate
is shortened block by block, so no appeal to an unbounded search is hidden
in the proof.
\end{proof}

\begin{theorem}[Right-linear grammars]\label{reg:grammar}
Finite right-linear grammars and finite automata define the same languages.
This remains true when productions have forms $A\to uB$ and $A\to u$,
where $u$ is any finite terminal word, including $\eps$.
\end{theorem}
\begin{proof}
First take productions $A\to aB$, $A\to B$, $A\to a$, and
$A\to\eps$.  Make one state for each nonterminal and one fresh accepting
state $f$.  Use the start nonterminal as initial state.  The four
production types give respectively an $a$ edge to $B$, an $\eps$ edge
to $B$, an $a$ edge to $f$, and an $\eps$ edge to $f$.  State $f$ has
no outgoing edges.  Before the last production, a derivable sentential
form is a terminal prefix followed by exactly one nonterminal.  Each
production extends the prefix and updates that state, while the final
production removes it.  This invariant converts a finite derivation
into an accepting path and conversely.
For a general terminal word $u$, replace the production by a private
chain spelling $u$; use an $\eps$ edge if $u$ is empty.  The same invariant
holds at the chain boundaries.
Conversely an NFA gives productions $A_q\to aA_r$ for every letter edge,
$A_q\to A_r$ for every $\eps$ edge, and $A_q\to\eps$ for every final
state.  A fresh start nonterminal has unit productions to initial states.
Again derivations and paths convert by scanning their finite steps.
All added chains and productions have explicit lengths supplied by the
input grammar or automaton, so the conversions are primitive recursive.
\end{proof}

\paragraph{Application of Arden's equation.}
The general PR-language result~\ref{alg:arden} applies in particular to
regular languages.  For $\eps\notin A$, the language $A^*B$ is the unique
PR solution of $X=AX\cup B$.  The closure constructions above additionally
produce its finite automaton from those for $A,B$.  This is useful when
solving finite systems of state equations; the explicit elimination
recurrence in Theorem~\ref{reg:eliminate} does not need to assume that
any particular coefficient is $\eps$-free.

\subsection{Pumping and explicit separation examples}

\begin{theorem}[Regular pumping with computed cuts]\label{reg:pumping}
Let $M$ be a DFA with $n$ states accepting $w$ with $|w|\geq n$.  Primitive
recursive functions compute $x,y,z$ such that
\[
 w=xyz,\qquad |xy|\leq n,\qquad 1\leq|y|,\qquad
 xy^tz\in L(M)\quad\hbox{for every }t\in\N.
\]
\end{theorem}
\begin{proof}
Inspect the first $n+1$ states of the run.  Finite pigeonhole supplies
$i<j\leq n$ with equal states; a bounded search selects the first such
pair.  Put $x=w[0:i]$, $y=w[i:j]$, $z=w[j:|w|]$.  The run on $y$ loops
at the state reached by $x$.  Quantifier-free induction on $t$ proves
that $y^t$ also loops there.  Appending the fixed suffix $z$ reaches the
same accepting state as the original run.  The cut and length assertions
follow from $i<j\leq n$.
\end{proof}

\begin{corollary}[Four explicit nonregular languages]\label{reg:nonregular}
Over letters $a\ne b$, the languages
\[
 \{a^rb^r:r\in\N\},\qquad
 \{w:w=w^R\},\qquad
 \{ww:w\in\{a,b\}^*\}
\]
are not regular.  The unary language $\{a^{r^2}:r\in\N\}$ is also not
regular.  Each assertion is proved in \ipra\ as the refutation of every
purported finite DFA code recognizing that particular PR predicate.
\end{corollary}
\begin{proof}
Suppose an $n$-state DFA recognizes the first language.  Apply
Theorem~\ref{reg:pumping} to $a^nb^n$.  The nonempty pumped factor lies
in the first $n$ letters and is therefore $a^d$ with $1\leq d\leq n$.
Taking $t=0$ accepts $a^{n-d}b^n$, whose two counts differ, a contradiction.
For palindromes intersect a hypothetical DFA with the regular language
$a^*ba^*$.  The intersection is exactly $\{a^rba^r:r\in\N\}$.
For its own number $n$ of states, pump $a^nba^n$ at the first $n$ letters
and take $t=0$; the two $a$-blocks no longer have equal lengths.
For squares of words intersect with $a^*ba^*b$.  A square in this slice
has exactly one $b$ in each half, and each half ends in $b$, so the
intersection is exactly $\{a^rba^rb:r\in\N\}$.  Pumping its initial
$a$-block gives the same contradiction.
Finally an $n$-state unary-square DFA accepts $a^{n^2}$.  Pumping with
$t=2$ gives an accepted length $n^2+d$, where $1\leq d\leq n$.
But $n^2<n^2+d<(n+1)^2$, so this length is not a square.  Deciding square
membership itself uses the bounded search $r\leq|w|$; all four predicates
and all contradiction witnesses are primitive recursive.
\end{proof}

\subsection{Finite decision procedures and canonical minimization}

\begin{theorem}[Emptiness and shortest accepted words]\label{reg:empty}
An $n$-state DFA has a nonempty language if and only if it accepts a word
of length below $n$.  Emptiness and a shortest accepted word, when one
exists, are primitive recursive.
\end{theorem}
\begin{proof}
An accepting run can be shortened by deleting loops, preserving its
initial and final states; the remaining path has fewer than $n$ edges.
Conversely any such word witnesses nonemptiness.  Enumerate words of
length $0,1,\ldots,n-1$, lexicographically within each length, and choose
the first accepted one.  The finite enumeration tests
$\sum_{i<n}k^i$ words.  If every test fails, no accepting run exists by
the shortening argument.
\end{proof}

\begin{theorem}[Equality, inclusion, universality, and short separators]
\label{reg:equality}
For DFAs with $n,m$ states, inequality of their languages has a witnessing
word of length below $nm$.  Failure of inclusion has the same bound.
Equality, inclusion, disjointness, and universality are primitive
recursive decisions on finite automaton codes.
\end{theorem}
\begin{proof}
Use the product automaton, accepting respectively when the two bits
differ, when the first is true and the second false, or when both are
true.  Lemma~\ref{reg:run} identifies its language with the required
counterexamples.  Theorem~\ref{reg:empty} bounds every shortest witness
by $nm-1$ and decides emptiness.  Universality is emptiness of the
complement.  A negative decision produces an explicit counterexample;
a positive decision says a finite Boolean reachability table has no
accepting state.  These are complete decisions, rather than searches
that stop only on counterexamples.
\end{proof}

\begin{theorem}[Finiteness has a bounded test]\label{reg:finiteness}
An $n$-state DFA accepts arbitrarily long words if and only if it accepts
a word of length in $[n,2n)$.  If it accepts no word in this interval,
all its accepted words have length below $n$.  Thus finiteness is decidable,
and a finite accepted language can be completely enumerated.
\end{theorem}
\begin{proof}
A word of length at least $n$ has a pumping decomposition.  Repeating its
nonempty factor gives words of unbounded length; for a requested bound
$B$, choosing $t=B+1$ already makes the pumped word longer than $B$.
Conversely, given any accepted word of length at least $n$, while its
length is at least $2n$, remove the nonempty loop furnished by its first
$n+1$ run states.  The removed length is at most $n$, so the new length
is still at least $n$.  Repeat at most the original length times.  Each
successful pass decreases the length, and the result lies in $[n,2n)$.
This is a bounded primitive recursive normalization of a supplied
accepted word.  If the finite test of that interval fails, no accepted
word can have length at least $n$.  Enumerating the shorter words then
lists the whole language.  If it succeeds, the pumping construction
produces an injection of $\N$ into the language, since the lengths
$|xz|+t|y|$ strictly increase.  This gives a constructive, witnessed
meaning of infinitude without using an infinite-pigeonhole principle.
\end{proof}

\begin{proposition}[Bounded distinguishability of states]\label{reg:distinguish}
For an $n$-state DFA, whether two states have identical future acceptance
is decidable by a primitive recursive table.  If they differ, a
separating suffix of length at most $n^2$ can be computed.
\end{proposition}
\begin{proof}
Put $D_0(q,r)=F(q)\mathbin{\mathrm{xor}}F(r)$ and iterate
\[
 D_{j+1}(q,r)=D_j(q,r)\vee
   \bigvee_{a<k}D_j(\delta(q,a),\delta(r,a)).
\]
The invariant is that $D_j(q,r)=1$ exactly when a word of length at most
$j$ distinguishes $q,r$.  At a newly true entry store the responsible
letter and predecessor entry; an initially true entry stores $\eps$.
The recursion either keeps a witness or prefixes a letter to one, proving
both directions by bounded induction.
The table is monotone and has only $n^2$ bits.  If a round makes no
change, every later round is identical.  If there were a change after
round $n^2$, every earlier round would have changed as well; this would
require more than $n^2$ newly set bits.  Finite counting therefore proves
$D_{n^2}=D_{n^2+1}$.  At this fixed point, a false pair has equal accepting
bits and sends every letter to a false pair.  Induction along any supplied
word now proves identical future acceptance.  A true pair has its
recorded suffix of length at most $n^2$.
\end{proof}

\begin{theorem}[Constructive DFA minimization]\label{reg:minimize}
There is a primitive recursive procedure that takes any DFA to an
equivalent reachable DFA with the smallest possible number of states.
It returns access words for its states and distinguishing suffixes for
every pair of different states.  All these witnesses have explicit
bounds in the size of the input automaton.
\end{theorem}
\begin{proof}
First restrict to states reachable from $q_0$, using
the bounded reachability result~\ref{base:graph}; choose a shortest access word of length below $n$
for each.  Write $q\sim r$ when $D_{n^2}(q,r)=0$.
The characterization in Proposition~\ref{reg:distinguish} makes this
an equivalence relation.  More explicitly, reflexivity and symmetry are
immediate from equality of accepting bits, and transitivity follows by
comparing the bits for every word of length at most $n^2$.  Its fixed-point
equation gives $q\sim r\Rightarrow\delta(q,a)\sim\delta(r,a)$.
Choose the least reachable representative of each equivalence class by
bounded search.  These representatives, ordered numerically and
renumbered consecutively, are the states of the quotient.  Transitions
send a representative to the representative of its successor; acceptance
is its original bit.  The right-congruence property makes the choice
independent of representatives, and induction on a run proves equality
of the accepted language.
To prove minimality, let $B$ be any other DFA for that language.  For
each quotient state choose the already computed access word $u_i$ and
map it to the state reached in $B$ by $u_i$.  For $i\ne j$ the table gives
a suffix $v_{ij}$ of length at most $n^2$ accepted from exactly one of
the two original states.  If $B$ reached the same state after $u_i,u_j$,
run composition would force the same answer on $u_iv_{ij}$ and
$u_jv_{ij}$, contrary to agreement with the original automaton.  Thus
the finite map is injective.  Finite pigeonhole implies that $B$ has at
least as many states as the quotient.  Every comparison used here is a
Boolean comparison of finite runs; the universally quantified assumption
that $B$ recognizes the same language is only instantiated at the
explicitly constructed words.
\end{proof}

\begin{corollary}[Uniqueness of the reachable minimal DFA]\label{reg:unique}
Two reachable minimal complete DFAs for the same language are isomorphic
by a unique initial-state-preserving, transition-preserving map.
The map and its inverse are obtained by bounded access-word searches.
\end{corollary}
\begin{proof}
Every state of a minimal DFA is distinguishable from every other: if
not, the quotient construction of Theorem~\ref{reg:minimize} would
strictly reduce its state count while preserving the language.
For a state $q$ of the first automaton choose a shortest access word $u_q$
and map $q$ to the state reached by $u_q$ in the second.  If $u,v$ reach
the same first state, the two corresponding second states accept the
same suffixes, because both automata recognize the same language.
Distinguishability in the second automaton forces those states to be
equal.  The map therefore does not depend on the chosen access word.
It preserves the initial state, transitions, and final bits.  The symmetric
construction is an inverse, as seen by using the same access words.
Finally any map preserving the initial state and transitions must send
the state reached by $u$ to the state reached by $u$; reachability makes
this condition determine the whole map.
\end{proof}

\begin{theorem}[Effective Myhill--Nerode certificate]\label{reg:nerode}
Fix a primitive recursive language predicate $L$.  A finite right-congruence
certificate consists of a fixed PR function $c$ with values in $[m]$,
words $r_0,\ldots,r_{m-1}$ with $c(r_i)=i$, and the assertions
\[
\begin{split}
 c(u)=c(v)&\Longrightarrow c(ua)=c(va),\qquad a<k,\\
 c(u)=c(v)&\Longrightarrow L(u)=L(v).
\end{split}
\]
Such a certificate gives an $m$-state DFA for $L$.  Conversely every DFA
for $L$ supplies such a certificate after removing inaccessible states.
Consequently regularity is equivalent to the existence of this effective
finite data, with their displayed arithmetic correctness assertions.
\end{theorem}
\begin{proof}
Use initial state $c(\eps)$, transition
$\delta(i,a)=c(r_i a)$, and accepting bit $F(i)=L(r_i)$.
If the current state after $u$ is $c(u)$, then $c(r_{c(u)})=c(u)$,
so the first certificate assertion gives
$\delta(c(u),a)=c(ua)$.  Induction along the word proves the run state
is $c(w)$; the second assertion gives $F(c(w))=L(w)$.
Conversely let $c(w)$ be the state reached after $w$ in a reachable DFA,
renumbering its states below $m$ if necessary.  The run-composition lemma
proves right stability and agreement of accepting bits.  Bounded
reachability provides a representative word for every state.
This is a schema for fixed PR functions, not an assertion that all PR
programs admit a PR universal evaluator.  Nor does it infer decidable
classes or representatives from an abstract finite-index equivalence
relation on arbitrary predicates.  Those additional data are exactly
what make the converse constructive and its realization primitive
recursive.
\end{proof}

\subsection{Lean implementation obligations for this section}
\label{reg:lean}

The semantic Lean layer may use \texttt{List (Fin k)} for words and the
standard \texttt{DFA}, \texttt{NFA}, and \texttt{\ensuremath{\eps}NFA}
interfaces.  It must not be confused with the arithmetic layer.  For the
latter use a record of natural-number fields for the alphabet size, state
count, start, transition-table code, and final-bit-vector code, together
with a PR Boolean predicate \texttt{validDFA}.  The NFA code has an initial
bit-vector and a three-dimensional Boolean edge table; its epsilon variant
adds a two-dimensional edge table.  Define a total decoder for every field,
using the rejecting automaton on invalid inputs.  A regular expression is
a postfix list of tagged nodes, each child index strictly less than its
parent index. The validity predicate also requires a single final root,
exactly one parent for every nonroot node, and connectivity to that root;
thus the code describes a tree, not a shared directed acyclic graph. These
conditions are bounded PR checks.  Grammar codes are lists of the
four production tags used in Theorem~\ref{reg:grammar}.

The following are the substantive proof obligations; none is silently
assumed from an algorithm's successful execution.
\begin{enumerate}
\item \texttt{runCode\_append}: decoded iteration equals list evaluation,
with the invariant ``the first $i$ scanned letters lead to the stored
state.''  Construct the code for this equality in the object calculus.
\item \texttt{shortenPath\_correct}: verify every deletion, preserve endpoints,
and combine the explicit finite-pigeonhole witness with the decreasing
length counter.  The output bound is $n-1$.
\item \texttt{epsilonClosure\_correct} and \texttt{subset\_correct}: prove the
closure table's witness invariant, then the subset vector's prefix
invariant.  The bounded epsilon-path bound is
$\ell+(\ell+1)(n-1)$.
\item \texttt{closure\_constructors\_correct}: each closure operator must
return a valid finite code and maps in both directions between the relevant
finite path certificates.  For quotients the crucial search bound is $nm$.
\item \texttt{regexTable\_correct}, \texttt{thompson\_correct}, and
\texttt{eliminate\_correct}: use the common finite table of expression
nodes, word intervals, and state endpoints.  For elimination also record
the internal-vertex threshold and the expression-size majorant.  Structural
induction in ambient Lean alone is insufficient evidence of an \ipra{}
proof; the table invariant is the route to quantifier-free induction.
\item \texttt{regularGrammar\_correct}: implement stepwise derivation/path
conversions and prove the terminal-prefix/nonterminal-suffix invariant.
\item \texttt{pump\_correct}: extract the repeated indices from the first
$n+1$ run states and prove the loop equation by quantifier-free induction
on the pumping exponent.  State the nonregularity results as explicit
refutations of alleged DFA recognition equations.
\item \texttt{equivalence\_decision\_correct},
\texttt{finiteness\_decision\_correct}, and
\texttt{minimize\_correct}: return either bounded witnesses or checked
negative reachability tables.  For minimization add the $n^2$-round
stabilization proof, representative correctness, and the finite injection
into any competing DFA.
\item \texttt{finiteRightCongruence\_toDFA}: treat the fixed PR code $c$ and
its object-theoretic correctness derivations as parameters.  Keep the
schema separate from the internal quantification over finite DFA data.
\end{enumerate}

For example, after defining the constructors the semantic target for
subset construction is
\[
 \forall M,w\quad
 \operatorname{validENFA}(M)=1\ \longrightarrow\
 \operatorname{acceptDFA}(\operatorname{determinize}(M),w)
 =\operatorname{acceptENFA}(M,w).
\]
Both acceptance terms here are primitive recursive: the second uses the
bounded path normalization above.  The object-level target is a derivation
of this formula in the encoded \ipra\ calculus, after its Boolean and
bounded quantifiers have been compiled away.  Existing mathlib declarations
such as \texttt{NFA.toDFA\_correct} and
\texttt{\ensuremath{\eps}NFA.toNFA\_correct} are useful semantic comparison
lemmas; they do not themselves discharge that target.  The inspected
regular-expression documentation explicitly describes its syntax and
semantics; no version-independent assumption about a preexisting complete
Kleene-equivalence implementation is needed here \cite{MathlibRegex}.
